\documentclass[11pt]{article}

\usepackage[a4paper,margin=27mm]{geometry}
\usepackage[T1]{fontenc}
\usepackage{lmodern}
\usepackage{amsmath,amssymb,amsthm,mathtools}
\usepackage{microtype}
\usepackage{xcolor}
\usepackage[colorlinks=true,linkcolor=blue!55!black,
  urlcolor=blue!55!black]{hyperref}
\usepackage[nameinlink,noabbrev]{cleveref}

\newtheorem{theorem}{Theorem}[section]
\newtheorem{lemma}[theorem]{Lemma}
\newtheorem{corollary}[theorem]{Corollary}
\theoremstyle{definition}
\newtheorem{definition}[theorem]{Definition}
\newtheorem{construction}[theorem]{Construction}
\theoremstyle{remark}
\newtheorem{remark}[theorem]{Remark}

\newcommand{\R}{\mathbb R}

\title{The no-four-concyclic variant V4:\\[1mm]a compact one-lemma proof}
\author{Companion note}
\date{5 August 2026}

\begin{document}
\maketitle

\begin{abstract}
For \(n\ge4\), consider \(n\) distinct real planar points which are not all
collinear and no four of which lie on one proper Euclidean circle.  Count
the distinct proper circles determined by selected noncollinear triples.
We prove that the minimum is
\[
 \binom{n-1}{2},
\]
and that equality occurs precisely for a near-pencil: \(n-1\) points on a
line and one point off it.  The proof identifies circles with noncollinear
triples and applies one elementary richest-line count.  It uses no finite
search, classification, induction, or external incidence theorem.
\end{abstract}

\section{The frozen variant and the result}

\begin{definition}[V4 configuration]\label{def:v4}
A \emph{V4 configuration} is a finite set \(P\subset\R^2\) of \(n\ge4\)
distinct points such that \(P\) is not collinear and no proper Euclidean
circle contains four points of \(P\).  A circle is counted if it contains a
noncollinear triple from \(P\); it is counted once, independently of its
full selected trace.  Write \(C(P)\) for this count and \(f_{\mathrm{V4}}(n)\)
for its minimum.
\end{definition}

The condition \emph{not collinear} is part of the definition.  Without it,
a collinear set would determine no proper circle and the minimum would be
zero.  Lines themselves are never counted.

\begin{definition}[Near-pencil]
A noncollinear \(n\)-point set is a \emph{near-pencil} if some line contains
exactly \(n-1\) of its points.
\end{definition}

\begin{theorem}\label{thm:v4}
For every integer \(n\ge4\),
\begin{equation}\label{eq:v4-value}
 \boxed{f_{\mathrm{V4}}(n)=\binom{n-1}{2}}.
\end{equation}
Moreover, a V4 configuration attains this value if and only if it is a
near-pencil.  In particular, the values for \(n=4,\ldots,10\) are
\[
 3,6,10,15,21,28,36.
\]
\end{theorem}

\section{The elementary circle interface}

We record the entire geometric input, so that the later extremal argument
is purely combinatorial.

\begin{lemma}\label{lem:circle-interface}
Three distinct noncollinear real points determine a unique proper Euclidean
circle.  A proper circle and a line have at most two common points.
\end{lemma}

\begin{proof}
A circle has an equation
\begin{equation}\label{eq:circle-equation}
 x^2+y^2+Ax+By+D=0.
\end{equation}
Substitution of three points gives a linear system for \(A,B,D\) whose
coefficient determinant is twice their signed triangle area.  It is
nonzero exactly when the points are noncollinear.  Thus the solution is
unique; because the equation contains three distinct real points, its
radius is positive.

After a rigid change of coordinates, a given line is \(y=0\).  Restricting
\eqref{eq:circle-equation} to it gives \(x^2+Ax+D=0\), which has at most two
real roots.  This proves the intersection bound.
\end{proof}

Let \(\tau(P)\) denote the number of noncollinear three-subsets of \(P\).

\begin{corollary}[Triple--circle bijection]\label{cor:triple-circle}
For every V4 configuration,
\begin{equation}\label{eq:circle-triple-bijection}
 C(P)=\tau(P).
\end{equation}
\end{corollary}

\begin{proof}
By \cref{lem:circle-interface}, every noncollinear triple determines one
proper circle.  If two distinct triples determined the same circle, their
union would contain at least four distinct selected points on that circle,
contrary to \cref{def:v4}.  The map is therefore injective.  It is
surjective because a counted circle is, by definition, one containing a
selected noncollinear triple.
\end{proof}

\section{One extremal lemma}

The next statement no longer mentions circles and is valid for every
finite noncollinear real point set.

\begin{lemma}[Richest-line triple bound]\label{lem:triple-bound}
Let \(P\subset\R^2\) be a set of \(n\ge4\) distinct points which are not all
collinear.  Then
\begin{equation}\label{eq:triple-bound}
 \tau(P)\ge\binom{n-1}{2}.
\end{equation}
Equality holds if and only if \(P\) is a near-pencil.
\end{lemma}

\begin{proof}
Choose a line \(L\) containing the largest possible number \(m\) of points
of \(P\), and put
\[
 k=n-m=|P\setminus L|.
\]
A line through a selected pair shows \(m\ge2\), while noncollinearity gives
\(k\ge1\).

First choose two points of \(P\cap L\) and one point outside \(L\).  These
are exactly
\[
 k\binom m2
\]
noncollinear triples.  Next fix an unordered pair \(x,y\in P\setminus L\).
The line \(xy\ne L\) contains at most one selected point of \(L\).  Hence at
least \(m-1\) points of \(P\cap L\) form a noncollinear triple with
\(x,y\).  Different outside pairs give different triples, so this supplies
at least
\[
 \binom k2(m-1)
\]
further triples.  The two families are disjoint, since their members contain
respectively two and one points of \(L\).  Triples wholly outside \(L\) may
be omitted.  Consequently
\begin{align}
 \tau(P)
 &\ge k\binom m2+\binom k2(m-1)\notag\\
 &=\frac{k(m-1)(n-1)}2.\label{eq:two-family-count}
\end{align}
The exact excess of the displayed lower bound over the target is
\begin{equation}\label{eq:factorized-gap}
 \frac{k(m-1)(n-1)}2-\binom{n-1}{2}
 =\frac{(n-1)(k-1)(m-2)}2\ge0.
\end{equation}
This proves \eqref{eq:triple-bound}.

Suppose equality holds.  The sandwich formed by
\eqref{eq:two-family-count}--\eqref{eq:factorized-gap} forces
\(k=1\) or \(m=2\).  In the first case \(L\) contains \(n-1\) points, so
\(P\) is a near-pencil.  In the second case maximality of \(L\) says that no
three points of \(P\) are collinear.  Every triple is then noncollinear, and
\[
 \tau(P)-\binom{n-1}{2}
 =\binom n3-\binom{n-1}{2}
 =\frac{(n-1)(n-2)(n-3)}6>0,
\]
contrary to equality for \(n\ge4\).  Conversely, in a near-pencil the
noncollinear triples are exactly the outsider together with a pair from
the long line, so there are \(\binom{n-1}{2}\).  This completes the equality
classification.
\end{proof}

\begin{remark}\label{rem:equality}
The richest-line choice is needed only for the equality classification;
the lower count itself works for any line determined by two selected
points.  For \(n\ge4\), the \((n-1)\)-point line of a near-pencil is unique:
two such distinct lines would contain at least \(2n-3>n\) selected points
in their union.  Thus the equality statement classifies the incidence
type, not a unique metric realization.
\end{remark}

\section{Proof and construction}

For a V4 configuration, \cref{cor:triple-circle,lem:triple-bound} give
immediately
\[
 C(P)=\tau(P)\ge\binom{n-1}{2}.
\]

\begin{construction}[Near-pencil extremizers]\label{con:near-pencil}
Take
\[
 (0,0),(1,0),\ldots,(n-2,0),\qquad u=(0,1).
\]
The set is not collinear.  By \cref{lem:circle-interface}, a proper circle
contains at most two selected points of the line \(y=0\); there is only one
selected point outside that line.  Thus every proper circle contains at
most three selected points, and the construction is V4-admissible.

Its noncollinear triples are precisely \(u\) together with a pair on
\(y=0\).  Therefore \cref{cor:triple-circle} gives exactly
\(\binom{n-1}{2}\) counted circles.
\end{construction}

The construction proves the upper bound in \eqref{eq:v4-value}, while the
preceding inequality proves the lower bound.  If a V4 configuration attains
equality, \cref{lem:triple-bound} makes it a near-pencil; conversely, the
same construction argument applies to every near-pencil.  This proves
\cref{thm:v4}.

\section{Scope and computational boundary}

The theorem concerns the frozen variant in \cref{def:v4}.  It is not the
historical unrestricted Erd\H{o}s problem: the no-four-concyclic condition
is an additional hypothesis.  It is also distinct from the no-three-
collinear variant, since V4 permits arbitrarily rich lines.

The same formula holds at \(n=3\) if noncollinearity is required: one
triangle determines one circle.  If instead no three points are allowed to
be collinear, then every triple in a V4 configuration determines a different
circle, so the count is \(\binom n3\).  Such configurations exist for every
finite \(n\): add points successively outside the finitely many lines and
circles forbidden by the points already chosen.

No computation is a premise of this proof.  The companion programs
\path{verify_v4_minimal_lemma.py} and
\path{verify_v4_minimal_lemma_redteam.py} check the polynomial identities
\eqref{eq:two-family-count}--\eqref{eq:factorized-gap} coefficientwise,
without looping over \(n\), line sizes, point sets, triples, circles, or
configurations.  They do not formalize the elementary geometric facts in
\cref{lem:circle-interface}.  Earlier finite-window and coordinate checks
are redundant archival tests, not proof dependencies.  This note is not a
Lean or Coq formalization.

\end{document}
